% Figure: multi-client backend design.
% Handshake commands never touch the backend, so a per-command lock lets the
% destination complete its handshake while the source is still connected.
\begin{tikzpicture}[
  font=\scriptsize,
  >=Latex,
  node distance=8mm,
  blk/.style={draw, rounded corners=2pt, align=center, inner sep=4pt,
              minimum height=9mm},
  ar/.style={->, thick},
]
  % clients
  \node[blk, fill=green!10, minimum width=26mm] (c0) at (0,2.6)
        {connection \#0\\\tiny (source VMM)};
  \node[blk, fill=green!10, minimum width=26mm] (c1) at (0,-0.2)
        {connection \#1\\\tiny (destination VMM)};

  % handles
  \node[blk, fill=orange!14, minimum width=30mm] (h0) at (4.6,2.6)
        {\texttt{SharedBackend}\\\tiny id = 0};
  \node[blk, fill=orange!14, minimum width=30mm] (h1) at (4.6,-0.2)
        {\texttt{SharedBackend}\\\tiny id = 1};

  % backend
  \node[blk, fill=blue!10, minimum width=32mm, minimum height=26mm] (be) at (9.4,1.2) {};
  \node[font=\scriptsize\bfseries] at (9.4,2.1) {\texttt{XhciBackend}};
  \node[font=\tiny, align=center] at (9.4,1.2)
        {xHCI state\\slots, endpoints\\event ring\\interrupt line\\DMA mappings};

  \draw[ar] (c0) -- node[above, font=\tiny]{handshake} (h0);
  \draw[ar] (c1) -- node[above, font=\tiny]{handshake} (h1);
  \draw[ar] (h0) -- node[above, font=\tiny, align=center]{lock per\\command} (be);
  \draw[ar] (h1) -- node[below, font=\tiny, align=center]{lock per\\command} (be);

  % invariant box
  \node[draw, rounded corners=2pt, fill=red!5, align=left, inner sep=4pt,
        font=\tiny, text width=62mm, anchor=north west] at (-0.2,-2.0)
  {\textbf{Invariants enforced by the wrapper}\\[2pt]
   \textbullet\ \emph{stale teardown:} only the connection that most recently
     registered interrupts may issue \texttt{SetIrqs} without fds or
     \texttt{DmaUnmap};\\[1pt]
   \textbullet\ \emph{interrupt hand-over:} registering a new line re-raises one
     interrupt so the guest re-examines the event ring;\\[1pt]
   \textbullet\ \emph{DMA idempotence:} re-mapping the same guest range replaces
     the old mapping instead of failing.};
\end{tikzpicture}
